\documentclass[11pt]{article}
\usepackage{amsmath, amssymb, amsthm}
\usepackage{geometry}
\usepackage{hyperref}
\usepackage{array}
\usepackage{booktabs}
\geometry{margin=1in}

\newtheorem{theorem}{Theorem}
\newtheorem{lemma}[theorem]{Lemma}
\newtheorem{proposition}[theorem]{Proposition}
\newtheorem{corollary}[theorem]{Corollary}
\theoremstyle{definition}
\newtheorem{definition}[theorem]{Definition}
\newtheorem{observation}[theorem]{Observation}
\newtheorem{remark}[theorem]{Remark}
\newtheorem{conjecture}[theorem]{Conjecture}

\newcommand{\Pq}{\Pi_q}
\newcommand{\PH}{P_H}
\newcommand{\Hq}{H_q(S_n)}
\newcommand{\Z}{\mathbb{Z}}
\newcommand{\Q}{\mathbb{Q}}
\newcommand{\tr}{\operatorname{tr}}

\title{Middle Factor of the Eigenvalue Rosetta Stone:\\
$(1+q)^{\lfloor n/2\rfloor}$ Divisibility of $\tr_{V_\lambda}(\Pi_q)$\\[4pt]
\large{Verification at $n=3,4,5$; partial factorization argument}}
\author{Clio}
\date{2026-04-23}

\begin{document}
\maketitle

\begin{abstract}
Let $\Pq = \prod_{j=1}^{\binom{n}{2}}(1+T_{i_j})$ be the staircase
Bott--Samelson element of the Iwahori--Hecke algebra $\Hq$ with
$T_i^2=(q-1)T_i+q$, and let $H=\langle s_1,s_3,s_5,\ldots\rangle\cong
(S_2)^{\lfloor n/2\rfloor}$ with Poincar\'e polynomial $(1+q)^{\lfloor
n/2\rfloor}$. We verify the middle factor of the \emph{Eigenvalue Rosetta
Stone} conjecture $e_\lambda(q)=q^{n(\lambda)}(1+q)^{\lfloor
n/2\rfloor}p_\lambda(q)$ at $n=3,4,5$: for every irreducible
$S_n$-representation $V_\lambda$ with $V_\lambda^H\ne 0$, the trace
polynomial $\tr_{V_\lambda}(\Pq)\in\Z[q]$ is divisible by
$(1+q)^{\lfloor n/2\rfloor}$ in $\Z[q]$, and each nonzero eigenvalue of
$\Pq|_{V_\lambda}$ (in $\overline{\Q(q)}$) is divisible by $(1+q)^{\lfloor
n/2\rfloor}$.

We also establish a partial algebraic identity: at $n=4,5$ the staircase
factorizes as $\Pq=A\cdot\PH\cdot B$ where $\PH=\prod_{i\text{ odd}}(1+T_i)$
is the $H$-symmetrizer. This puts $(1+q)^{\lfloor n/2\rfloor}=
(1+q)^{\operatorname{rank}(\PH^2\to\PH)}$ into the trace
via $\PH^2=(1+q)^{\lfloor n/2\rfloor}\PH$, \emph{if} integrality of the
quotient trace can be shown. We document honestly where the argument
passes and where it stalls.
\end{abstract}

\section{Setup}

Let $\Hq$ be the Iwahori--Hecke algebra of $S_n$ with generators
$T_1,\ldots,T_{n-1}$ and relations
\[
  T_i^2=(q-1)T_i+q,\quad T_iT_{i+1}T_i=T_{i+1}T_iT_{i+1},\quad
  T_iT_j=T_jT_i\ (|i-j|\ge 2).
\]
Fix the \emph{staircase reduced word} for $w_0\in S_n$:
\[
  \mathbf{i}=(s_1)(s_2s_1)(s_3s_2s_1)\cdots(s_{n-1}\cdots s_1),\qquad
  \ell(\mathbf{i})=\tbinom{n}{2}.
\]
Define
\begin{equation}
  \Pq \;=\; \prod_{j=1}^{\binom{n}{2}}\bigl(1+T_{i_j}\bigr)
  \;\in\;\Hq,\qquad
  \PH \;=\; \prod_{i\in\{1,3,5,\ldots\}\cap[1,n-1]}(1+T_i).
\end{equation}
Write $s=\lfloor n/2\rfloor$, so $H=(S_2)^{s}$ is the parabolic with odd
generators and $|H|=2^s$, Poincar\'e polynomial
$\sum_{w\in H}q^{\ell(w)}=(1+q)^s$.

\medskip

Two facts used throughout:

\begin{lemma}\label{lem:idempotency}
$(1+T_i)^2=(1+q)(1+T_i)$; hence $\PH^2=(1+q)^s \PH$.
\end{lemma}

\begin{proof}
$(1+T_i)^2=1+2T_i+T_i^2=1+2T_i+(q-1)T_i+q=(1+q)(1+T_i)$. Since the odd
generators pairwise commute, $\PH=\prod_{i\text{ odd}}(1+T_i)$ has
$\PH^2=\prod_i(1+T_i)^2=(1+q)^s\PH$.
\end{proof}

\begin{lemma}\label{lem:Hinv}
For every irrep $V_\lambda$, $\Pq\cdot V_\lambda\subseteq V_\lambda^H$; in
fact $\operatorname{im}\Pq|_{V_\lambda}=V_\lambda^H$ when the latter is
nonzero. (``H-invariant theorem,'' proved 2026-04-17.)
\end{lemma}

\section{Main theorem}

\begin{theorem}[Middle factor, verified at $n\le 5$]\label{thm:main}
For $n\in\{3,4,5\}$ and every partition $\lambda\vdash n$ with
$V_\lambda^H\ne 0$:
\begin{enumerate}
\item[(a)] \textbf{Trace divisibility.}
  $\tr_{V_\lambda}(\Pq)$ is divisible by $(1+q)^s$ in $\Z[q]$, where
  $s=\lfloor n/2\rfloor$.
\item[(b)] \textbf{Per-eigenvalue divisibility.}
  Each of the $\dim V_\lambda^H$ nonzero eigenvalues of $\Pq|_{V_\lambda}$,
  as an element of $\overline{\Q(q)}$, is divisible by $(1+q)^s$ in the
  sense that its quotient by $(1+q)^s$ is integral over $\Z[q]$.
\end{enumerate}
\end{theorem}

\begin{proof}[Proof by direct computation.]
We use Young's seminormal form in the Hoefsmit/Ram convention. For each
standard Young tableau $T$ of shape $\lambda$ label pairs $(T,s_iT)$ with
axial distance $d=\mathrm{ct}(T[i])-\mathrm{ct}(T[i+1])$ and set
$a(d)=(q-1)/(1-q^d)$. On the 2-dimensional block spanned by $(v_T,v_{s_iT})$
with $T$ ordered before $s_iT$:
\[
  T_i v_T = a(d)\,v_T + v_{s_iT},\qquad
  T_iv_{s_iT} = \bigl(a(d)a(-d)+q\bigr)v_T + a(-d)\,v_{s_iT}.
\]
Same-row cells give $T_iv_T=qv_T$; same-column cells give $T_iv_T=-v_T$.
We verified for each $(\lambda,n)$ below that the matrices satisfy
$T_i^2=(q-1)T_i+q$, the braid relation, and commutativity for
$|i-j|\ge 2$. Then $\Pq|_{V_\lambda}$ is computed by matrix product along
the staircase word, and $\tr_{V_\lambda}(\Pq)$ is extracted as a
$\Q(q)$-expression and checked for polynomial division by $(1+q)^s$.

The results are exhibited in Table~\ref{tab:traces}. Every row with
$V_\lambda^H\ne 0$ has $\tr_{V_\lambda}(\Pq)\in(1+q)^s\,\Z[q]$; the
explicit quotient $\tr_{V_\lambda}(\Pq)/(1+q)^s$ is shown.

For (b): when $\dim V_\lambda^H=1$, the unique nonzero eigenvalue equals
the trace, so (a) already implies (b). When $\dim V_\lambda^H=2$ (occurs
only for $\lambda=(4,1),(3,2)$ at $n=5$), the two nonzero eigenvalues
$\alpha_1,\alpha_2$ satisfy $x^2-Sx+P=0$ with $S=\alpha_1+\alpha_2$ the
trace and $P=\alpha_1\alpha_2$ the subleading coefficient of the
characteristic polynomial. Substituting $y=x/(1+q)^s$ gives
$y^2-\bigl(S/(1+q)^s\bigr)y+\bigl(P/(1+q)^{2s}\bigr)=0$, which has
coefficients in $\Z[q]$ iff $(1+q)^s\mid S$ and $(1+q)^{2s}\mid P$. Both
are verified:
\begin{center}\small
\begin{tabular}{lll}
\toprule
$\lambda,n$ & $S=\tr\,\Pq$ & $P=\alpha_1\alpha_2$\\
\midrule
$(4,1),n=5$ & $q(1+q)^4(q^4+7q^3+17q^2+7q+1)$ &
  $q^3(1+q)^{12}(q^2+4q+1)$\\
$(3,2),n=5$ & $q^2(1+q)^2(q^4+9q^3+19q^2+9q+1)$ &
  $q^5(1+q)^8(q^2+4q+1)$\\
\bottomrule
\end{tabular}
\end{center}
In both cases $(1+q)^s$ divides $S$ (with $s=2$, divides to higher order in
fact) and $(1+q)^{2s}$ divides $P$.
\end{proof}

\begin{table}[hp]
\centering
\caption{Traces $\tr_{V_\lambda}(\Pq)$ for $n=3,4,5$, computed via
Hoefsmit seminormal form. $s=\lfloor n/2\rfloor$. Rows with
$V_\lambda^H=0$ have $\Pq|_{V_\lambda}=0$ by Lemma~\ref{lem:Hinv} and are
omitted.}
\label{tab:traces}
\small
\begin{tabular}{llllll}
\toprule
$n$ & $\lambda$ & $\dim V_\lambda$ & $\dim V_\lambda^H$ &
  $\tr_{V_\lambda}(\Pq)$ & $\tr_{V_\lambda}(\Pq)/(1+q)^s$ \\
\midrule
$3$ & $(3)$     & $1$ & $1$ & $(1+q)^3$ & $(1+q)^2$ \\
$3$ & $(2,1)$   & $2$ & $1$ & $q(1+q)$ & $q$ \\
\midrule
$4$ & $(4)$     & $1$ & $1$ & $(1+q)^6$ & $(1+q)^4$ \\
$4$ & $(3,1)$   & $3$ & $1$ & $q(1+q)^2(q^2+4q+1)$ & $q(q^2+4q+1)$ \\
$4$ & $(2,2)$   & $2$ & $1$ & $q^2(1+q)^2$ & $q^2$ \\
\midrule
$5$ & $(5)$     & $1$ & $1$ & $(1+q)^{10}$ & $(1+q)^8$ \\
$5$ & $(4,1)$   & $4$ & $2$ & $q(1+q)^4(q^4+7q^3+17q^2+7q+1)$
                & $q(1+q)^2(q^4+7q^3+17q^2+7q+1)$ \\
$5$ & $(3,2)$   & $5$ & $2$ & $q^2(1+q)^2(q^4+9q^3+19q^2+9q+1)$
                & $q^2(q^4+9q^3+19q^2+9q+1)$ \\
$5$ & $(3,1,1)$ & $6$ & $1$ & $q^3(1+q)^2(q^2+4q+1)$ & $q^3(q^2+4q+1)$ \\
$5$ & $(2,2,1)$ & $5$ & $1$ & $q^4(1+q)^2$ & $q^4$ \\
\bottomrule
\end{tabular}
\end{table}

\section{Partial algebraic argument via $\Pq=A\cdot\PH\cdot B$}

For $n=4,5$ the staircase word begins with four factors that align with
$\PH$ after a single commutation:

\begin{lemma}\label{lem:APHB}
For $n=4$: $\Pq=A\,\PH\,B$ with
\[
  A=(1+T_1)(1+T_2),\quad
  \PH=(1+T_1)(1+T_3),\quad
  B=(1+T_2)(1+T_1).
\]
For $n=5$:
\[
  A=(1+T_1)(1+T_2),\ \
  \PH=(1+T_1)(1+T_3),\ \
  B=(1+T_2)(1+T_1)(1+T_4)(1+T_3)(1+T_2)(1+T_1).
\]
\end{lemma}

\begin{proof}
The staircase word for $n=4$ is $s_1\cdot(s_2s_1)\cdot(s_3s_2s_1)$ of
length $6$. Read left-to-right:
$(1+T_1)(1+T_2)\,(1+T_1)(1+T_3)\,(1+T_2)(1+T_1)$. Positions $3,4$ give
generators $s_1,s_3$ which commute; their product is precisely $\PH$.
Identification of $A$ and $B$ by inspection. The $n=5$ case is the same
with $B$ extended by the first four factors of the $s_4\cdots s_1$ block.

Verified computationally in the length-lex regular representation of
$H_q(S_4)$ ($24\times 24$) and $H_q(S_5)$ ($120\times 120$): every entry
of $\Pq-A\PH B$ is identically zero.
\end{proof}

\begin{corollary}[Cyclic-trace reformulation]\label{cor:cyclic}
For $n=4,5$ and every representation $V$ of $\Hq$,
\[
  \tr_V(\Pq) = \tr_V(A\,\PH\,B) = \tr_V(\PH\cdot Y),\qquad Y:=BA\in\Hq.
\]
\end{corollary}

\begin{proof}
Cyclic property of trace.
\end{proof}

\subsection{What would close the proof algebraically}

Set $e_H:=\PH/(1+q)^s$. On any $\Hq$-module $V$ over $\Q(q)$,
$e_H$ is an idempotent (Lemma~\ref{lem:idempotency}) with
$\operatorname{im}e_H=V^H$. Corollary~\ref{cor:cyclic} combined with
Lemma~\ref{lem:idempotency} gives
\begin{equation}\label{eq:quot-trace}
  \tr_{V_\lambda}(\Pq)
  \;=\;(1+q)^s\,\tr_{V_\lambda}(e_H\cdot Y)
  \;=\;(1+q)^s\,\tr_{V_\lambda^H}(Y|_{V_\lambda^H}),
\end{equation}
where the last equality uses that $e_H$ is the projector onto
$V_\lambda^H\subseteq V_\lambda$ and cyclic trace of a projector.

\begin{proposition}[Trace divisibility, assuming integrality]\label{prop:conditional}
If $\tr_{V_\lambda^H}\bigl(Y|_{V_\lambda^H}\bigr)\in\Z[q]$ then
$(1+q)^s\mid\tr_{V_\lambda}(\Pq)$ in $\Z[q]$.
\end{proposition}

\begin{proof}
Direct from \eqref{eq:quot-trace}.
\end{proof}

The \emph{a priori} statement is only that this trace is an element of
$\Q(q)$. Integrality (no $(1+q)$ in the denominator) is the missing step.

\subsection{Where the naive strengthening fails}

One is tempted to argue: $\PH\cdot\Hq\cdot\PH\subseteq(1+q)^s\Hq$, whence
$\PH Y\PH$ is divisible in $\Hq$ and its trace inherits the divisibility.
This \emph{fails}:

\begin{proposition}[Counterexample at $n=4$]\label{prop:counter}
The ideal $\PH\Hq\PH\subseteq H_q(S_4)$ is not contained in
$(1+q)^2 H_q(S_4)$. Specifically, for the double coset
$Hs_2H\subseteq S_4$ one computes
$\PH\cdot T_{s_2}\cdot\PH\notin(1+q)^2 H_q(S_4)$; some coefficients in the
$T$-basis are in $\Z[q]\setminus(1+q)^2\Z[q]$.
\end{proposition}

\begin{proof}
Direct computation in the left-regular representation of $H_q(S_4)$.
Details in the verification script \texttt{check\_PH\_yP\_H\_v2.py};
the minimum $k$ such that $(1+q)^k$ divides every $T$-basis coefficient
of $\PH T_{s_2}\PH$ is $k=0$, not $k=2$.
\end{proof}

\begin{remark}
The double-coset decomposition $S_n=\coprod_d H\sigma_d H$ partitions
$\PH\Hq\PH=\sum_d\PH T_{\sigma_d}\PH$. For $d=e$ (trivial coset) and for
any coset where $\sigma_d$ has a reduced expression with a left- and a
right-descent in $H$, we have
$\PH T_{\sigma_d}\PH\in(1+q)^{2s}\Hq$ by the
$(1+T_i)T_{s_i}=q(1+T_i)$ pushing. The obstruction is cosets whose
minimal-length representative has \emph{no} descent in $H$ on either side
(the ``diagonal'' cosets), of which $\sigma=s_2$ at $n=4$ is the
smallest example.

This explains why the direct algebraic approach fails: the trace on any
single module can still be divisible when summed over cosets, because
spiky entries cancel --- but this cancellation is a representation-theoretic
fact ($V^H$ has low dimension), not an algebra-level fact.
\end{remark}

\subsection{Why Lemma~\ref{lem:APHB} does not extend to $n\ge 6$}

The staircase word for $w_0\in S_n$ reads
\[
  s_1\,|\,s_2s_1\,|\,s_3s_2s_1\,|\,s_4s_3s_2s_1\,|\,\cdots
\]
For the identity $\Pq=A\PH B$ we need a contiguous sub-word whose
indices, in some order, are precisely $\{1,3,5,\ldots\}\cap[1,n-1]$.
At $n=4$: positions $3,4$ carry $(s_1,s_3)$ --- exact match. At $n=5$:
same positions, same match. At $n=6$: $H=\langle s_1,s_3,s_5\rangle$
needs indices $\{1,3,5\}$ but the staircase has no contiguous trio of
commuting odd-indexed factors; the earliest occurrence of $s_5$ is
position $11$, and the three $s_1,s_3,s_5$ do not line up as a
factor. Thus the Lemma~\ref{lem:APHB} strategy does not extend.

\section{Verification summary}

\begin{itemize}
\item \textbf{Hoefsmit seminormal form.} Implemented per-irrep as
  $d\times d$ symbolic matrices over $\Q(q)$ with $d=\dim V_\lambda$.
  Checked: $T_i^2=(q-1)T_i+q$, braid, far-commutativity for all irreps at
  $n=3,4,5$. ({\tt extract\_n5\_eigenvalues.py})
\item \textbf{Trace divisibility table.} All entries of
  Table~\ref{tab:traces} extracted via
  $\Q(q)$-matrix products then symbolic polynomial division.
  Ten irreps verified; every one passes.
\item \textbf{Factorization $\Pq=A\PH B$.} Confirmed in the
  $n!\times n!$ length-lex regular representation: all $576$ ($n=4$) and
  $14{,}400$ ($n=5$) entries of $\Pq-A\PH B$ are zero.
  ({\tt verify\_APHB.py})
\item \textbf{Per-eigenvalue divisibility.} From characteristic
  polynomials of $\Pq|_{V_\lambda}$: for $\dim V_\lambda^H\ge 2$ cases
  $(4,1)$ and $(3,2)$ at $n=5$, verified $(1+q)^s\mid S$ and
  $(1+q)^{2s}\mid P$ as above.
\item \textbf{Counterexample for the algebra-level claim.}
  Confirmed that $\PH T_{s_2}\PH\notin(1+q)^2 H_q(S_4)$ by exhibiting
  specific coefficients. ({\tt check\_PH\_yP\_H\_v2.py})
\end{itemize}

\section{What remains}

\begin{enumerate}
\item \textbf{Integrality of $\tr_{V_\lambda^H}(Y|_{V_\lambda^H})$
  (primary gap).} Proposition~\ref{prop:conditional} reduces trace
  divisibility to proving that the trace on $V_\lambda^H$ of $BA$
  (acting via $\Hq$) lies in $\Z[q]$, not merely in
  $(1+q)^{-?}\Z[q]$. This is equivalent to showing that the idempotent
  $e_H=\PH/(1+q)^s$ has an integer lift on $V_\lambda$ --- i.e., the
  projector onto $V_\lambda^H$ is representable by an element of
  $\Z[q]_{(1+q)}\cdot\Hq$ with well-defined reduction $\mod(1+q)$. Over
  the DVR $\Z[q]_{(1+q)}$, $\PH$ itself is nilpotent modulo $(1+q)$, so
  no naive lift exists; a more refined argument (parabolic Hecke
  algebra, base change to the residue field) is needed.
\item \textbf{Extension to $n\ge 6$.} The
  $\Pq=A\PH B$ factorization does not persist; the staircase word no
  longer admits the contiguous $\{1,3,5,\ldots\}$ subword. A different
  organization --- e.g.\ the Soergel/EW-category
  $\Pq=BS(\mathbf{i})=\bigoplus_w m_w B_w$ --- would be insensitive to
  this accident of word position.
\item \textbf{Palindromic factor $p_\lambda(q)$.} The last column of
  Table~\ref{tab:traces} exhibits, after dividing by $q^{n(\lambda)}$,
  the proposed palindromic factors: for $n=5$,
  $p_{(5)}=(1+q)^8$, $p_{(4,1)}=(1+q)^2(q^4+7q^3+17q^2+7q+1)$,
  $p_{(3,2)}=q^4+9q^3+19q^2+9q+1$, $p_{(3,1,1)}=q^2+4q+1$,
  $p_{(2,2,1)}=1$. Each is palindromic and has non-negative integer
  coefficients. Interpreting $p_\lambda$ as a graded trace or KL
  polynomial remains open.
\end{enumerate}

\section*{Honest scope of this note}

Theorem~\ref{thm:main} is a \emph{computational} theorem at $n\le 5$, not
a general-$n$ structural proof. The partial algebraic skeleton
(Lemma~\ref{lem:APHB}, Corollary~\ref{cor:cyclic},
Proposition~\ref{prop:conditional}) reduces the question to an
integrality statement we do not yet know how to prove, and in any case
fails to extend to $n\ge 6$. The Soergel/BS($\mathbf{i}$) route --- where
$(1+q)^s$ is the Poincar\'e polynomial of the parabolic $H$ and
$V^H$-multiplicities are graded --- remains the strongest candidate for
a full proof. This note records the precise state of the proof effort as
of 2026-04-23.

\bigskip

\noindent\emph{Code and data:} the seminormal-form implementation and
trace/divisibility verification scripts are in the scratch directory and
reproduce Table~\ref{tab:traces} and every claim labelled
``computational verification'' above.

\end{document}
